\documentclass[11pt]{article}
\usepackage[a4paper,margin=21mm]{geometry}
\usepackage[no-math]{fontspec}
\setmainfont{Latin Modern Roman}
\newfontfamily\CJKfont{Noto Serif CJK SC}
\XeTeXlinebreaklocale "zh"
\XeTeXlinebreakskip=0pt plus 1pt
\usepackage{amsmath,amssymb,amsthm,amscd,graphicx,color}
\usepackage{xurl}
\usepackage[unicode,hidelinks]{hyperref}
\allowdisplaybreaks
\setlength\emergencystretch{3em}
\numberwithin{equation}{section}

\numberwithin{equation}{section}

\newtheorem{thm}{Theorem}[section]
\newtheorem*{thm*}{Theorem}
\newtheorem{cor}[thm]{Corollary}
\newtheorem{lem}[thm]{Lemma}
\newtheorem{prop}[thm]{Proposition}
{ \theoremstyle{definition}
	\newtheorem{dfn}[thm]{Definition}
	\newtheorem{Note}[thm]{Note}
	\newtheorem{eg}[thm]{Example}
	\newtheorem{rmk}[thm]{Remark} }

\usepackage{tikz,tikz-cd}
\usepackage{slashed}
\usepackage{mathrsfs}

\def\cA{\mathcal A}\def\cB{\mathcal B}\def\cC{\mathcal C}\def\cD{\mathcal D}
\def\cE{\mathcal E}\def\cF{\mathcal F}\def\cG{\mathcal G}\def\cH{\mathcal H}
\def\cI{\mathcal I}\def\cJ{\mathcal J}\def\cK{\mathcal K}\def\cL{\mathcal L}
\def\cM{\mathcal M}\def\cN{\mathcal N}\def\cO{\mathcal O}\def\cP{\mathcal P}
\def\cQ{\mathcal Q}\def\cR{\mathcal R}\def\cS{\mathcal S}\def\cT{\mathcal T}
\def\cU{\mathcal U}\def\cV{\mathcal V}\def\cW{\mathcal W}\def\cX{\mathcal X}
\def\cY{\mathcal Y}\def\cZ{\mathcal Z}

\def\AA{\mathbb A}\def\BB{\mathbb B}\def\CC{\mathbb C}\def\DD{\mathbb D}
\def\EE{\mathbb E}\def\FF{\mathbb F}\def\GG{\mathbb G}\def\HH{\mathbb H}
\def\II{\mathbb I}\def\JJ{\mathbb J}\def\KK{\mathbb K}\def\LL{\mathbb L}
\def\MM{\mathbb M}\def\NN{\mathbb N}\def\OO{\mathbb O}\def\PP{\mathbb P}
\def\QQ{\mathbb Q}\def\RR{\mathbb R}\def\SS{\mathbb S}\def\TT{\mathbb T}
\def\UU{\mathbb U}\def\VV{\mathbb V}\def\WW{\mathbb W}\def\XX{\mathbb X}
\def\YY{\mathbb Y}\def\ZZ{\mathbb Z}

\def\sA{\mathscr A}\def\sB{\mathscr B}\def\sC{\mathscr C}\def\sD{\mathscr D}
\def\sE{\mathscr E}\def\sF{\mathscr F}\def\sG{\mathscr G}\def\sH{\mathscr H}
\def\sI{\mathscr I}\def\sJ{\mathscr J}\def\sK{\mathscr K}\def\sL{\mathscr L}
\def\sM{\mathscr M}\def\sN{\mathscr N}\def\sO{\mathscr O}\def\sP{\mathscr P}
\def\sQ{\mathscr Q}\def\sR{\mathscr R}\def\sS{\mathscr S}\def\sT{\mathscr T}
\def\sU{\mathscr U}\def\sV{\mathscr V}\def\sW{\mathscr W}\def\sX{\mathscr X}
\def\sY{\mathscr Y}\def\sZ{\mathscr Z}

\def\fa{\mathfrak a}\def\fb{\mathfrak b}\def\fc{\mathfrak c}\def\fd{\mathfrak d}
\def\fe{\mathfrak e}\def\ff{\mathfrak f}\def\fg{\mathfrak g}\def\fh{\mathfrak h} \def\fj{\mathfrak j}\def\fk{\mathfrak k}\def\fl{\mathfrak l} \def\fm{\mathfrak m}\def\fn{\mathfrak n}\def\fo{\mathfrak o}\def\fp{\mathfrak p} \def\fq{\mathfrak q}\def\fr{\mathfrak r}\def\fs{\mathfrak s}\def\ft{\mathfrak t}
\def\fu{\mathfrak u}\def\fv{\mathfrak v}\def\fw{\mathfrak w}\def\fx{\mathfrak x}
\def\fy{\mathfrak y}\def\fz{\mathfrak z}

\def\cat#1{\ensuremath{\mathsf{#1}}}
\DeclareMathOperator{\Coh}{\cat{Coh}}
\DeclareMathOperator{\QCoh}{\cat{QCoh}}
\DeclareMathOperator{\Mod}{\cat{Mod}}
\DeclareMathOperator{\Mult}{\cat{Mult}}

\def\lie#1{\ensuremath{\mathfrak{#1}}}
\def\aut{\ensuremath{\lie{aut}}}
\def\aff{\ensuremath{\lie{aff}}}
\def\so{\lie{so}}

\def\bu{\bullet}

\def\d{{\rm d}}

\DeclareMathOperator{\Spec}{Spec}
\DeclareMathOperator{\Spin}{Spin}
\DeclareMathOperator{\End}{End}
\DeclareMathOperator{\Sym}{Sym}
\DeclareMathOperator{\id}{id}
\DeclareMathOperator{\Hom}{Hom}
\DeclareMathOperator{\Vect}{Vect}
\DeclareMathOperator{\Ext}{Ext}

\def\sL{\mathscr L}
\def\sR{\mathscr R}

\makeatletter\newlabel{sec: examples}{{5}{}{Original source locator}{}{}}\makeatother
\begin{document}
\begin{center}\Large Derived pure-spinor dg equivalence for strict-object multiplets of two-step super-Poincare-type algebras\end{center}
\noindent\textbf{Stable ID:} 1091\quad\textbf{Author:} Chris Elliott; Fabian Hahner; Ingmar Saberi\par
\noindent\textbf{Work:} The Derived Pure Spinor Formalism as an Equivalence of Categories\par
\noindent\textbf{Source:} \url{https://sigma-journal.com/2023/022/}\par
\noindent\textbf{Version:} Official SIGMA 19 (2023) article 022, DOI 10.3842/SIGMA.2023.022; journal PDF and native TeX acquired 2026-09-30, exact bytes pinned by SHA256\par
\noindent\textbf{Rights:} CC BY 4.0. Authors retain copyright. \url{https://creativecommons.org/licenses/by/4.0/}. Reuse and typesetting adaptation; original language and mathematical content preserved.\par
\noindent\textbf{Difficulty (prior selection preserved):} {\CJKfont H3: The author proof constructs the two inverse functors and realizes both composites by explicit de Rham and Koszul contractions. It gives finite homotopy-perturbation formulas, then computes the transferred higher module operations to recover the original supersymmetry action. Naturality on morphisms closes the equivalence of dg-categories, rather than merely identifying cohomology groups or invoking abstract Koszul duality.}\par
\medskip\noindent\textbf{Editorial boundary note (not author text):} Complete original Theorem4.3, both explicit inverse functors and full naturality-on-morphisms proof. All differential/total parity and affine-geometric action conventions, strict and homotopy multiplet definitions and strict-object category, two-step superalgebra and full automorphism hypotheses, finite-degree equivariant Chevalley–Eilenberg modules, left/right supertranslation actions, derived pure-spinor functor and current de Rham/Koszul/finite component-field constructions are retained. The complete current homotopy-data and finite perturbation formulas, transferred higher module operations and recovery of the original action close both composite functors. Named formal homotopy-transfer/strictification and Hilbert-resolution facts remain the exact author references, without recursive reproof. One expressly named principal dg-equivalence outcome; examples, antifield applications and later consequences are unselected, and supporting constructions receive no separate count. All author maps, shifts, finite-degree conditions and literal formulas are preserved.\par
\noindent\textbf{External original reference sec: examples:} Original applications/examples section, source1135–1452. Mentioned in grading and component-field exposition only; no selected theorem proof step is deferred to it.\par
\medskip\hrule\medskip\noindent\textbf{Author text begins below.}\par
\setcounter{section}{1}
% Original source lines 255--452
\section{The category of multiplets} \label{sec: multiplets}
In~\cite{perspectives}, the notion of a \emph{$\fg$-multiplet} for a super Lie algebra $\fg$ was formalized. Let us quickly review the relevant definitions, and then move on to define morphisms between $\fg$-multiplets and study the resulting category $\Mult_{\fg}$. We start with the case of strict multiplets and strict morphisms and then discuss the homotopy-theoretic generalization.

In brief, by a \emph{multiplet} on a smooth manifold $X$ we will mean a graded vector bundle $E$ on $X$ together with a differential $D$ and an action $\rho$ of an appropriate super Lie algebra $\fg$, compatible with an action of the even part of $\fg$ on the base manifold $X$. The reader should consider the prototypical situation where $X = \RR^n$, and $\fg$ is a super Lie algebra whose even part is the Poincar\'e algebra of infinitesimal isometries of $\RR^n$. A multiplet can be thought of as modelling the perturbations of a classical field in the BV formalism.

\begin{rmk}
	We would like to clearly emphasize here the distinction between multiplets and theories in this context. In general, the data of a multiplet is \emph{less data} than the data of a classical field theory: it does not include enough data to specify dynamics. The data of a perturbative classical supersymmetric field theory in the BV formalism includes the data of a multiplet, together with some additional data. If a multiplet is equipped with the additional data of a~\emph{BV bracket} (an odd shifted symplectic structure), it can be viewed as modelling a \emph{free} theory in the BV formalism. Such structures can easily be constructed from multiplets by taking the sum of a multiplet with a shift of its dual.	If one wishes to include non-trivial interactions, one must additionally equip the sections of the bundle $E$ with an $L_\infty$ structure, compatibly with the supersymmetry action.
	
	So, while the results of this paper do apply to BV field theories, when we speak of \emph{multiplets} we are referring to a less constrained notion, and to obtain a genuine field theory from a multiplet one needs to specify strictly more data.
\end{rmk}

\subsection{Gradings}
\label{ssec:gradings}
Many objects appearing throughout this work (e.g., vector spaces, vector bundles, Lie algebras etc.) will carry a grading by $\ZZ\times \ZZ/2\ZZ$ as well as a differential of bidegree $(1,+)$. We refer to such objects with the prefix ``dgs'' standing for ``differential graded super''. Typically these gradings have a clear physical meaning: the integer grading refers to the ghost number or cohomological degree, while the $\ZZ/2\ZZ$-grading corresponds to the intrinsic parity (fermion number modulo two). The total parity denotes the $\ZZ/2\ZZ$-grading which arises by forgetting the $\ZZ$-grading to~$\ZZ/2\ZZ$ and then totalizing with the intrinsic parity. It is the total parity which governs all signs.

It will often be convenient to introduce an additional piece of data: an auxiliary action of a one-dimensional abelian Lie algebra $\RR$ on a dgs vector space. Many of our examples will be built from the following simple observation.

\begin{eg}
	Let $\fn$ be a dg Lie algebra concentrated in degrees 1 and 2. For all non-negative integers $k$, there is a natural action of $\RR$ on $\Sym^k(\fn^*)$, where $\RR$ acts on duals to the degree 1 and 2 summands of $\fn^*$ with weight $-1$ and $-2$ respectively.
\end{eg}

The results in this paper will make sense for arbitrary $\RR$-equivariant dgs vector spaces, but our most common examples will have the following behavior.

\begin{dfn}
	A \emph{lifted} dgs vector space is an $\RR$-equivariant dgs vector space where the $\RR$-action has integral weights, and the $\ZZ/2\ZZ$-grading coincides with the $\RR$-weight modulo two.
\end{dfn}

Let us now introduce some notation that we will use when we discuss lifted dgs vector bundles on a smooth manifold $X$. Let us write
\begin{equation*}
E = \bigoplus_{(w,d) \in \ZZ^2} E^{w,d}
\end{equation*}
for a lifted dgs vector bundle, where the first index indicates the decomposition of $E$ by $\RR$-weight, and the second index indicates internal $\ZZ$-degree on $E$. In some of our calculations, particular in~Section~\ref{sec: examples}, we will place the summands in a two-dimensional array, where the two coordinates are given by $w-d$ and $d$. For example,
\begin{equation*}
E =
\begin{bmatrix}
\cdots &E^{-1,0} &E^{0,0} &E^{1,0} &E^{2,0}&\cdots\\
\cdots &E^{0,1} &E^{1,1} &E^{2,1} &E^{3,1}&\cdots\\
\cdots &E^{1,2} &E^{2,2} &E^{3,2} &E^{4,2} &\cdots
\end{bmatrix}.
\end{equation*}
So we recover $w$ as the total degree with respect to this sheared bigrading. The overall parity is then determined by the column.

If the differential on our dgs vector bundle decomposes as a sum of homogeneous differential operators of order $k$, say $D = \sum_{k \ge 0} D_k$, the summand $D_k$ acts---with respect to the sheared grading---with bidegree $(2k-1, 1)$. So the homogeneous summands all increase the vertical degree by 1, but they may modify the horizontal degree by any odd integer $\ge -1$.

\subsection{Strict multiplets}

For a dgs vector bundle\footnote{Note that each graded piece $E^k$ is a finite rank vector bundle, but the total rank of $E$ may still be infinite.} $(E,D)$ over a base smooth manifold $X$, we denote the space of global smooth sections by $\cE = \Gamma(X,E)$. The endomorphisms $\End(\cE)$ form a dgs Lie algebra, where the bracket is given by the commutator and the differential is $[D,-]$. Inside $(\End(\cE), [D,-])$ there is a sub dgs Lie algebra consisting of all endomorphisms which act on sections via differential operators. We denote this subalgebra by $(\cD(E), [D,-])$.
\begin{dfn}
	A \emph{strict local dgs $\fg$-module} is a triple $(E,D,\rho)$ where $(E,D)$ is a dgs vector bundle and
	\begin{equation*}
	\rho \colon \ \fg \longrightarrow (\cD(E) , [D,-] )
	\end{equation*}
	is a map of dgs Lie algebras. Here the super Lie algebra $\fg$ is viewed as a dgs Lie algebra in cohomological degree zero with trivial differential.
\end{dfn}
\begin{rmk}
	This definition (as well as many of the following) has a natural generalization for~$\fg$ a dgs Lie algebra. Since we are ultimately interested in the pure spinor superfield formalism, we restrict our attention to super Lie algebras with no cohomological grading.
\end{rmk}
Note that, since a super Lie algebra $\fg$ has no differential, $\rho$ commutes with the differential on the dgs vector bundle,
\begin{equation*}
[D,\rho(x)] = 0 \qquad \forall x \in \fg .
\end{equation*}
It is standard to encode a $\fg$-module structure $\rho$ on $(E,D)$ as a dgs Lie algebra structure on the direct sum $\fg \oplus \cE$. Concretely we set for the unary and binary operations
\begin{align}
&\mu_1(x_1,\sigma_1) = (0, D\sigma_1),\nonumber \\
&\mu_2((x_1,\sigma_1),(x_2,\sigma_2)) = \big([x_1,x_2] , \rho(x_1)\sigma_2 - (-1)^{|\sigma_1| |x_2|} \rho(x_2)\sigma_1 \big) ,\label{eq: lie structure on sum}
\end{align}
where $x_1,x_2 \in \fg$ and $\sigma_1,\sigma_2 \in \cE$.

There is an obvious notion of morphisms between strict local $\fg$-modules.
\begin{dfn}
	A \emph{strict morphism} of strict local $\fg$-modules	$(E,D,\rho)$ and $(E',D',\rho')$ is a map of cochain complexes
	\begin{equation*}
	\psi \colon \ \cE \longrightarrow \cE'
	\end{equation*}
	realized by differential operators such that
	\begin{equation*}
	\psi \circ \rho(x) = \rho'(x) \circ \psi
	\end{equation*}
	for all $x \in \fg$.
\end{dfn}
A strict morphism $\psi \colon (E,D,\rho) \longrightarrow (E',D',\rho')$ gives rise to a strict morphism of the associated dgs Lie algebras by setting $\tilde{\psi} = \id_{\fg} \times \psi \colon \fg \oplus \cE \longrightarrow \fg \oplus \cE'$. Conversely it is easy to check that every strict morphism of dgs Lie algebras of that form gives rise to a strict morphism of $\fg$-modules. We call $\psi$ a quasi-isomorphism if it is a quasi-isomorphism of dgs vector bundles; equivalently $\tilde{\psi}$ is a quasi-isomorphism of dgs Lie algebras.

Now, suppose that $(E,D)$ is a dgs vector bundle over an affine space $V$ (so, in the notation used above, $V=X$). Let us additionally assume that $\fg$ is equipped with a map of super Lie algebras
\begin{equation*}
\phi \colon \ \aff(V) \longrightarrow \fg .
\end{equation*}
In essence, when we refer to strict $\fg$-\emph{multiplets}, we are referring to strict local $\fg$-modules for which the affine transformations acts in a geometric way.
\begin{dfn}
	A \emph{strict $\fg$-multiplet} $(E,D,\rho)$ on $V$ is an affine dgs vector bundle over $V$ equipped with a strict local $\fg$-module structure
	\begin{equation*}
	\rho \colon \ \fg \longrightarrow \cD(E)
	\end{equation*}
	such that the pullback of the module structure along~$\phi$ agrees with the natural action on sections of the affine vector bundle. Concretely, this means that the following diagram commutes:
	\begin{equation} \label{eq: geom action}
	\begin{tikzcd}
	\fg \arrow[r, "\rho"] & \cD(E). \\
	\aff(V) \arrow[ru, "\mathrm{aff}" '] \arrow[u,"\phi"]
	\end{tikzcd}
	\end{equation}
\end{dfn}
Here $\mathrm{aff}$ denotes the natural action of the affine algebra on $E$.
We define the category of strict multiplets with strict morphisms to be the full subcategory of the category of strict local $\fg$-modules with objects strict $\fg$-multiplets. We denote this category by $\Mult_{\fg}^{\mathrm{strict}}$.
%
\begin{rmk}
	The condition~\eqref{eq: geom action} stating that the translations act geometrically heavily depends on $\phi$. In particular, since the action $\mathrm{aff}$ is effective, the corresponding category of $\fg$-multiplets is trivial if $\phi$ is not injective.
\end{rmk}
\begin{rmk}
	Note that a morphism of strict local $\fg$-modules is automatically compatible with the action of $\aff(V)$. For example let $\psi \colon (E,D,\rho) \longrightarrow (E',D',\rho')$ be a strict morphism. Then we have
	\begin{equation*}
	\psi \circ \mathrm{aff}(x) = \psi \circ \rho(\phi(x)) = \rho'(\phi'(x)) \circ \psi = \mathrm{aff}'(x) \circ \psi .
	\end{equation*}
	Therefore it is sensible to define $\Mult_\fg^{\text{strict}}$ as a full subcategory of strict local $\fg$-modules.
\end{rmk}
%
\subsection{Homotopy theory of multiplets}
%
It is well known that various strict algebraic structures, such as associative algebras or Lie algebras, admit homotopy-theoretic generalizations ($A_\infty$ or $L_\infty$ algebras in these cases). The same is true for strict $\fg$-module structures,
which generalize to $L_\infty$ $\fg$-modules. In the context of this work, we will refer to such homotopy module structures simply as module structures and choose to emphasize whenever a module is strict instead. In this spirit we can easily define (not necessarily strict) local $\fg$-modules and $\fg$-multiplets by replacing Lie maps in the above definitions with $L_\infty$ maps.

Thus, a local $\fg$-module is just a dgs vector bundle $(E,D)$ together with an $L_\infty$ map of $L_\infty$ algebras
\begin{equation*}
\rho \colon \ \fg \rightsquigarrow \cD(E) .
\end{equation*}
Recall that this means there are component maps
\begin{equation*}
\rho^{(k)} \colon \ \fg^{\otimes k} \longrightarrow \cD(E)[1-k], \qquad k \geq 1
\end{equation*}
satisfying a series of compatibility relations the lowest of which reads
\begin{equation*}% \label{eq: L-action}
\big[\rho^{(1)}(x),\rho^{(1)}(y)\big] - \rho^{(1)}([x,y]) =\big[D, \rho^{(2)}(x,y)\big] \qquad \forall x,y \in \fg.
\end{equation*}
Clearly a local $\fg$-module is strict if and only if $\rho^{(k)} = 0$ for all $k\geq2$.

\begin{rmk}
	As we discussed in the introduction, considering this homotopical version of a~module for a super Lie algebra is very natural in supersymmetric field theory. In most supersymmetric classical field theories, when one describes the most natural multiplet of fields, there is only a strict action of the supersymmetry algebra on-shell, i.e., after imposing the equations of motion. If one, however, allows $L_\infty$ actions, it is possible to describe an $L_\infty$ action of the supersymmetry algebra on the space of fields in the BV formalism directly, without introducing auxiliary fields. This idea goes back to the beginnings of the BV formalism; applications to global symmetries are discussed, for example, in~\cite{BaulieuSusy}. The example of super Yang--Mills theory is discussed rigorously and systematically in all dimensions in~\cite{ESW}.
\end{rmk}


Similarly to the strict case, we can conveniently encode a $\fg$-module structure $\rho$ as an~$L_\infty$ structure on $\fg \oplus \cE$. To this end we supplement the operations~\eqref{eq: lie structure on sum} by the following brackets for $k\geq3$ (for details we refer to~\cite{Allocca, Lada})
\begin{equation*}
\mu_k((x_1,v_1),\dots,(x_k,v_k)) = \left( 0 , \sum_{i=1}^{k} \pm \rho^{(k-1)}\big(x_1,\dots,\hat{x}_i,\dots,x_k\big)v_i \right).
\end{equation*}
One can define a morphism $\psi \colon (E,D,\rho) \longrightarrow (E',D',\rho')$ of local $\fg$-modules by component maps
\begin{equation*}
\psi_n \colon \ \fg^{\otimes n-1} \otimes \cE \longrightarrow \cE' ,
\end{equation*}
which are given by differential operators and satisfy a series of compatibility relations (see~\cite{Allocca} for details). We recover strict morphisms by restricting to those where the only component map is $\psi_1$. Again, we can describe such morphisms by morphisms of the associated $L_\infty$ algebras
\begin{equation*}
\tilde{\psi} \colon \ \fg \oplus \cE \rightsquigarrow \fg \oplus \cE'
\end{equation*}
by supplementing $\tilde{\psi}_1 = \id_{\fg} \times \psi_1$ with
\begin{equation*}
\tilde{\psi}_k ((x_1,\sigma_1) , \dots , (x_k, \sigma_k)) = \left(0, \sum_{i=1}^{k} \pm \psi_k\big(x_1, \dots , \hat{x}_i , \dots , x_k , \sigma_i \big)\right)
\end{equation*}
for $k \geq 2$. $\psi$ is a quasi-isomorphism of local $\fg$-modules if $\psi_1$ is a quasi-isomorphism of cochain complexes, or equivalently if $\tilde{\psi}$ is a quasi-isomorphism of $L_\infty$ algebras. In general, encoding module structures as $L_\infty$ structures is very convenient because it allows the use of many known tools from the theory of $L_\infty$ algebras, like homotopy transfer for instance.

\begin{rmk}
	An equivalent realization of the notion of a non-strict morphism between a pair of local $L_\infty$ algebras $\fg,\fg'$ is provided by considering the Chevalley--Eilenberg chain complex~$C_\bullet(\fg)$ of an $L_\infty$ algebra $\fg$. This is a cocommutative dg coalgebra whose underlying graded coalgebra is $\Sym^\bullet(\fg[-1])$, with differential induced from the $L_\infty$ structure on $\fg$ as a sum of terms given by the $L_\infty$ brackets on $\fg$. The data of a morphism $\psi \colon \fg \to \fg'$ of local $L_\infty$ algebras is equivalent to that of a morphism of cocommutative dg coalgebras
	\[\psi' \colon \ C_\bullet(\fg) \to C_\bullet(\fg')\]
	given by differential operators. This interpretation allows us to view local $L_\infty$ algebras as objects of a dg-category, so that we can discuss (for example) homotopies between $L_\infty$ algebra morphisms. When the source $\fg$ is finite-dimensional, we can dually consider morphisms of commutative dg algebras $C^\bullet(\fg') \to C^\bullet(\fg)$. This is for example the case for the local $\fg$-module structures $\rho \colon \fg \rightsquigarrow \cD(E)$ appearing in the definition of multiplets.
\end{rmk}

Let us now give the definition of a (not necessarily strict) multiplet.
\begin{dfn}
	A \emph{$\fg$-multiplet} $(E,D,\rho)$ is an affine dgs vector bundle over $V$ equipped with a local $\fg$-module structure
	\begin{equation*}
	\rho\colon \ \fg \rightsquigarrow \cD(E)
	\end{equation*}
	such that the pullback of the module structure along~$\phi$ agrees strictly with the natural action on sections of the affine vector bundle. Concretely, this means that the following diagram commutes:
	\begin{equation*}
	\begin{tikzcd}
	\fg \arrow[r, "\rho^{(1)}"] & \cD(E). \\
	\aff(V) \arrow[ru, "\mathrm{aff}" '] \arrow[u,"\phi"]
	\end{tikzcd}
	\end{equation*}
	We define the dg-category of $\fg$-multiplets to be the full subcategory of local $\fg$-modules with objects being $\fg$-multiplets and denote it by $\Mult_{\fg}$.
\end{dfn}

\begin{rmk}
	Note that, in particular, the action $\rho \circ \phi$ obtained by restricting the $L_\infty$ action to affine transformations is necessarily strict.
\end{rmk}

We will sometimes wish to refer to the full subcategory of strict multiplets, but allowing arbitrary morphisms.
\begin{dfn}
	Denote by $\Mult_{\fg}^{\text{\rm strict-ob}}$ the full sub dg-category of $\Mult_{\fg}$ generated by strict multiplets.
\end{dfn}


We can obtain homotopy categories from the dg-categories $\Mult_{\fg}$ as well as $\Mult_{\fg}^{\mathrm{strict}}$ by replacing the hom space $\Hom_{\Mult_{\fg}}(E, E')$ by its zeroth cohomology $\mathrm H^0(\mathrm{Hom}_{\Mult_{\fg}}(E, E'))$. We denoted the resulting categories by $\mathrm{Ho}(\Mult_{\fg})$ and $\mathrm{Ho}(\Mult_{\fg}^{\mathrm{strict}})$. Speaking physically, quasi-isomorphisms of $\fg$-multiplets correspond to \emph{perturbative} equivalences of multiplets, where by ``perturbative'' here we mean equivalences of the derived formal neighborhoods of a point in the classical moduli field space. Therefore isomorphism classes in the homotopy category correspond to perturbatively distinct multiplets.


\begin{rmk}
	Because every $L_\infty$ algebra can be strictified \cite[Part II, Corollary 1.6]{KrizMay}, the natural inclusion
	\begin{equation*}
	\mathrm{Ho}\big(\Mult_{\fg}^{\text{\rm strict-ob}}\big) \to \mathrm{Ho}(\Mult_{\fg})
	\end{equation*}
	is an equivalence of categories, where $\Mult_{\fg}^{\text{\rm strict-ob}}$ denotes the full subcategory of multiplets spanned by strict objects.
\end{rmk}

\setcounter{section}{2}
% Original source lines 519--1008
\section{Derived pure spinor superfields}
In this section, we define the derived generalization of the pure spinor superfield construction, which will provide one of the two functors that witness the equivalence of categories. We begin by setting up the general context in which we want to work.

As in the introduction, let $\fn$ be a two-step nilpotent super Lie algebra, defined by a central extension
\begin{equation}\label{eq:2step}
	0 \to \fn_2 \to \fn \to \Pi \fn_1 \to 0
\end{equation}
of the odd abelian super Lie algebra $\Pi \fn_1$. We imagine such objects as generalizations of supertranslation algebras. The automorphisms of $\fn$, which are all outer, will contain an abelian factor generating scale transformations, with respect to which $\fn_1$ has weight one and $\fn_2$ weight two. There is also a natural map
\begin{equation*}
	\aut(\fn) \to \lie{gl}(\fn_2).
\end{equation*}
The kernel of this map can be thought of as the $R$-symmetry algebra; in physical examples, $\aut(\fn)$ will be the product of $\so(\fn)$, scale transformations, and the $R$-symmetry algebra.

All of our constructions will take place in reference to a fixed super Lie algebra $\fg$ of the following type:
\begin{dfn}
	A super Lie algebra $\fg$ is of \emph{super Poincar\'e type} if it can be written as an extension
	\begin{equation*}
	0 \to \fn \to \fg \to \fg_0 \to 0,
	\end{equation*}
	where $\fn$ is a two-step nilpotent super Lie algebra of the form~\eqref{eq:2step} and $\fg_0$ is a Lie algebra equipped with a Lie map $\fg_0 \to \aut(\fn)$.
\end{dfn}

This means that the $\ZZ/2\ZZ$ grading on $\fg$ can be lifted to a $\ZZ$-grading concentrated in degrees zero, one, and two:
\begin{equation*}
\fg = \fg_0 \oplus \fn_1 \oplus \fn_2.
\end{equation*}
The most important examples are the super Poincar\'e algebras in various dimensions and with various amounts of supersymmetry.

It is possible to consider examples where $\fg_0$ is any subalgebra of~$\aut(\fn)$. We will typically wish, however, to take $\fg_0$ to
be such that any linear transformation of $\fn_2$ arising from an automorphism of~$\fn$ can be generated by an element of~$\fg_0$.
In the super-Poincar\'e case, this will mean that $\so(n)$ is contained in~$\fg_0$. In this paper, we will always simply take $\fg_0 = \aut(\fn)$.

Let us choose a basis $d_\alpha$ for $\fn_1$ and $e_\mu$ for $\fn_2$ such that we can expand the symmetric bracket in structure constants\footnote{In the context of super Poincar\'e algebras, these structure constants are typically expressed in terms of the matrix elements of the gamma matrices.}
\begin{equation*}
\big[d_\alpha , d_\beta \big] = f^\mu_{\alpha \beta} e_\mu .
\end{equation*}
We denote by $R = \Sym^\bullet\big(\fn_1^*\big) = \CC[\lambda^\alpha]$ the ring of polynomial functions on $\fn_1$. For $Q$ in $\fn_1$, the equation $[Q,Q] = 0$ defines an ideal $I$ in $R$, explicitly given by
\begin{equation*}
I = \big(\lambda^\alpha f^\mu_{\alpha \beta} \lambda^\beta \big) .
\end{equation*}
As we will discuss shortly, we can identify the quotient algebra $R/I$ with the degree zero Lie algebra cohomology $\mathrm H^0(\fn)$. Here, by degree zero, we mean with respect to the totalization of the canonical $\ZZ \times \ZZ$-grading on Chevalley--Eilenberg cochains.

The pure spinor superfield formalism gives a systematic tool to construct $\fg$-multiplets from the input datum of a graded $\fg_0$-equivariant $R/I$-module. Here we generalize the pure spinor superfield construction to $\fg_0$-equivariant $C^\bullet(\fn)$-modules and show that it defines a functor.


\subsection[The category of C\^{}bullet(n)-modules]{The category of $\boldsymbol{C^\bullet(\fn)}$-modules} \label{CE_module_section}
Recall that the Chevalley--Eilenberg complex of $\fn$ takes the form
\begin{equation*}
C^\bullet(\fn) = \big( \Sym^\bullet\big(\fn^*[1]\big) , \d_{\mathrm{CE}} \big),
\end{equation*}
where the Chevalley--Eilenberg differential $\d_{\mathrm{CE}}$ is induced by the dual of the bracket. Here the notation $\fn^*[1]$ means that we shift $\fn^*$ down in cohomological degree by one. The Chevalley--Eilenberg complex has a $\ZZ \times \ZZ/2\ZZ$-grading endowing it with the structure of a dgs algebra. In our present case, since the the $\ZZ/2\ZZ$ grading of $\fn$ lifts to a $\ZZ$-grading where $\fn_i$ has weight $i$, $C^\bullet(\fn)$ can be given a $\ZZ \times \ZZ$-grading. Totalizing this grading, the generators of $\fn_1^*$ sit in degree $0$ while the generators of $\fn_2^*$ live in degree $-1$. We denote these generators by $\lambda^\alpha$ and $v^\mu$ respectively. We can thus identify{\samepage
\begin{equation*}
C^{-p}(\fn) = \wedge^p \fn_2^* \otimes R
\end{equation*}
with respect to the totalized grading.}

The Chevalley--Eilenberg differential acts on these generators according to
\begin{equation*}
\begin{split}
\d_{\mathrm{CE}} \lambda^\alpha &= 0, \\
\d_{\mathrm{CE}} v^\mu &= \lambda^\alpha f^\mu_{\alpha \beta} \lambda^\beta.
\end{split}
\end{equation*}
In coordinates we will often write the Chevalley--Eilenberg algebra in the form
\begin{equation*}
\left( C^\bullet(\fn), \d_{\mathrm{CE}} \right) = \left( \CC[\lambda^\alpha ,v^\mu] , \d_{\mathrm{CE}} = \lambda^\alpha f^\mu_{\alpha \beta} \lambda^\beta \frac{\partial}{\partial v^\mu} \right) .
\end{equation*}
Using this description, we immediately see that we indeed recover $H^0(\fn) = R/I$.

$C^\bullet(\fn)$ is a dgs algebra, therefore we can consider dgs modules over it.
\begin{dfn}
	A $C^\bullet(\fn)$-module is a dgs vector space $(\Gamma, \d_\Gamma)$ together with a morphism
	\begin{equation*}
	(C^\bullet(\fn) , \d_{\mathrm{CE}}) \longrightarrow (\End(\Gamma) , [\d_\Gamma , -]),
	\end{equation*}
	of dgs algebras.
\end{dfn}



\begin{dfn}
	A morphism of $C^\bullet(\fn)$-modules $(\Gamma , \d_\Gamma)$ and $(\Gamma' , \d_{\Gamma'})$ is a cochain map
	\begin{equation*}
	f \colon \ (\Gamma , \d_\Gamma) \longrightarrow (\Gamma' , \d_{\Gamma'})
	\end{equation*}
	such that
	\begin{equation*}
	f(x \cdot_\Gamma \gamma) = x \cdot_{\Gamma'} f(\gamma) .
	\end{equation*}
	In other words, $f$ is just a morphism of dgs modules.
\end{dfn}
Note that $\fg_0$ acts on both $\fn_1$ and $\fn_2$ and thus also on $C^\bullet(\fn)$.

\begin{dfn}
	A $C^\bullet(\fn)$ module $\Gamma$ is called \emph{$\fg_0$-equivariant} if $\Gamma$ is also a representation of $\fg_0$ and the module structure map is $\fg_0$-equivariant. We further assume that each degree $\Gamma^k$ is finite dimensional as a complex vector space. We will denote the dg-category of $\fg_0$-equivariant $C^\bullet(\fn)$-modules by $\Mod_{C^\bullet(\fn)}^{\fg_0}$.
\end{dfn}

Our main results will take place in this category of $\fg_0$-equivariant modules.

\subsection{The derived pure spinor superfield formalism}

There is a canonical strict multiplet associated to the super Lie algebra $\fg$: the free superfield. Recall that $\fn$ is a two-step nilpotent super Lie algebra and let $N = \exp(\fn)$ be the associated nilpotent super Lie group. Since all such nilpotent super Lie groups are split, $N$ can be chosen to be the total space of a bundle with fiber $\Pi N_1$ over $N_2$, where $N_2$ is a vector space treated as an additive abelian group. (As super vector spaces, $\fn$ and~$N$ are isomorphic; the group operation can be constructed using the Baker--Campbell--Hausdorff formula, which terminates at quadratic order in this case.)

Left and right translations induce two commuting, $\aut(\fn)$-equivariant actions of $\fn$ on $N$ by vector fields
\begin{equation*}
\sL,\, \sR \colon \ \fn \longrightarrow \Vect(N) .
\end{equation*}
Let us choose coordinates $x^\mu$ on the abelian group $N_2$
and $\theta^\alpha$ on $N_1$. In terms of these coordinates, the vector fields given by the action of the odd elements can be described as follows:
\begin{equation*}
\begin{split}
\sR(d_\alpha) &= \frac{\partial}{\partial \theta^\alpha} - f^\mu_{\alpha \beta} \theta^\beta \frac{\partial}{\partial x^\mu}, \\
\sL(d_\alpha) &= \frac{\partial}{\partial \theta^\alpha} + f^\mu_{\alpha \beta} \theta^\beta \frac{\partial}{\partial x^\mu}.
\end{split}
\end{equation*}
In addition, the even elements simply act by derivatives
\begin{equation*}
\sR(e_\mu) = \sL(e_\mu) = \frac{\partial}{\partial x^\mu}.
\end{equation*}
The free superfield is the strict $\fg$-multiplet with $\cE = C^\infty(N)$, vanishing differential, and module structure given by the left translations $\sL$. Note that by ``free'' here, we mean in the sense of the rank one free module over the ring of smooth functions on superspace (``unconstrained''), rather than ``non-interacting''. Recall that all our multiplets are a priori non-interacting: the inclusion of interactions is additional data with which a multiplet can be equipped.


Let us now generalize the pure spinor superfield formalism to a derived setting.
\begin{dfn}
	We define the \emph{pure spinor functor} $A^\bullet \colon \Mod_{C^\bullet(\fn)}^{\fg_0} \longrightarrow \Mult_{\fg}^{\mathrm{strict}}$ by setting
	\begin{equation*}
	A^\bullet(\Gamma) = \big(C^\infty(N) \otimes_\CC \Gamma , \cD\big) ,
	\end{equation*}
	for an object $(\Gamma , \d_\Gamma)$.
	The differential $\cD$ is constructed using the right action $R$ and the $C^\bullet(\fn)$-module structure on $\Gamma$. Explicitly, it takes the form
	\begin{equation*}
	\cD = \lambda^\alpha \sR(d_\alpha) + v^\mu \sR(e_\mu) + \d_\Gamma .
	\end{equation*}
	For a morphism $f\colon\Gamma \longrightarrow \Gamma'$ we define
	\begin{equation*}
	A^\bullet(f) = \id_{C^\infty(N)} \otimes f \colon \ A^\bullet(\Gamma) \longrightarrow A^\bullet(\Gamma') .
	\end{equation*}
\end{dfn}

A few comments are in order.

\begin{enumerate}\itemsep=0pt
	\item The differential $\cD$ squares to zero precisely since $\Gamma$ is a $C^\bullet(\fn)$-module.
	\item We can equip $A^\bullet(\Gamma)$ with the structure of a dgs vector bundle over the spacetime $V := N_2$ by placing $C^\infty(N)$ in cohomological degree zero. Note that in particular the differential is of bidegree $(1,+)$.
	
	\item Further, $\fg$ acts on $C^\infty(N)$ via left translations and on $\Gamma$ by the trivial extension of the $\fg_0$-module structure which was part of the input datum. The tensor product of these two action makes $A^\bullet(\Gamma)$ into a strict multiplet.
	\item It is immediate to check that $A^\bullet(f)$ is a strict morphism of strict multiplets. Since $f$ is a~morphism of $C^\bullet(\fn)$-modules we have
	\begin{equation*}
	A^\bullet(f) \circ \cD = \cD' \circ A^\bullet(f) .
	\end{equation*}
	In addition,
	\begin{equation*}
	A^\bullet(f) \circ \sL = \sL \circ A^\bullet(f)
	\end{equation*}
	obviously follows from the definition.
	\item $A^\bullet$ is additive. The direct sum of two $C^\bullet(\fn)$-modules is mapped to the direct sum of the respective multiplets, $A^\bullet(\Gamma \oplus \Gamma') = A^\bullet(\Gamma) \oplus A^\bullet(\Gamma')$.
\end{enumerate}

This construction is a direct generalization of the pure spinor superfield formalism as described in~\cite{perspectives}. To see this, we notice that the category of graded equivariant of $R/I$-modules sits as a subcategory inside $\Mod_{C^\bullet(\fn)}^{\mathrm{strict}}$, namely precisely as those modules concentrated in cohomological degree zero. Indeed, every $R/I$-module is a $C^\bullet(\fn)$-module by the map
\begin{equation*}
C^\bullet(\fn) = \CC\big[\lambda^\alpha,v^\mu\big] \longrightarrow R/I, \qquad \big(\lambda^\alpha , v^\mu\big) \mapsto \lambda^\alpha .
\end{equation*}
The other way round, let $(\Gamma , \d_\Gamma)$ be a $C^\bullet(\fn)$-module concentrated in cohomological degree zero. Then the differential $\d_\Gamma$ vanishes and $v^\mu$ acts trivially for degree reasons. Therefore one has, for~$\gamma$ any element of $\Gamma$,
\begin{equation*}
0 = \d_\Gamma \big(v^\mu \cdot \gamma\big) = \big(\d_{\mathrm{CE}} v^\mu\big) \gamma = \big(\lambda f^\mu \lambda\big) \gamma ,
\end{equation*}
which endows $\Gamma$ with the structure of an $R/I$-module.

Restricting the functor $A^\bullet$ to graded equivariant $R/I$-modules, we obtain a functor
\begin{equation*}
A^\bullet_{R/I} \colon \ \Mod_{R/I}^{\fg_0} \longrightarrow \Mult_{\fg}^{\mathrm{strict}}
\end{equation*}
where the output simplifies to
\begin{equation*}
A^\bullet_{R/I} (\Gamma) = \left( C^\infty(N) \otimes_\CC \Gamma , \lambda^\alpha \sR(d_\alpha) \right)
\end{equation*}
such that we precisely recover the pure spinor superfield formalism as presented in~\cite{perspectives}.

Further, this is a derived generalization of the pure spinor superfield construction in the sense that $C^\bullet(\fn)$ can be viewed as a derived replacement of the ring $R/I$. We will therefore refer to this construction as the derived pure spinor superfield formalism.

\begin{rmk}
	We can alternatively view this geometrically as a derived enhancement of the affine nilpotence variety $\mathrm{Spec}(R/I)$. We can view a multiplet as arising from a quasi-coherent sheaf over the nilpotence variety, but this point of view requires forgetting some of the data given by the $C^\bullet(\fn)$ action. The philosophy of derived algebraic geometry suggests instead retaining this information by viewing a multiplet as arising from a coherent sheaf over the affine derived scheme $\mathrm{Spec}(C^\bullet(\fn))$---the derived analogue of the nilpotence variety.
\end{rmk}

\begin{rmk}
	In many cases the multiplet $A^\bullet(\Gamma)$ can be equipped with additional structures such as higher brackets yielding an $L_\infty$ structure. For example if $\Gamma$ is not only a $C^\bullet(\fn)$-module but in addition an algebra, one can take a Lie algebra $\fh$ such that the tensor product $\Gamma \otimes \fh$ forms a dgs Lie algebra. Then $A^\bullet(\Gamma \otimes \fh) = A^\bullet(\Gamma) \otimes \fh$ is also a dgs Lie algebra. An $L_\infty$ structure on the component fields arises via homotopy transfer. This was studied in the case of ten-dimensional super Yang--Mills theory in~\cite{AKLL, perspectives}.
\end{rmk}

\subsection[The multiplet associated to C\^{}bullet(n)]{The multiplet associated to $\boldsymbol{C^\bullet(\fn)}$} \label{sec: forms}

As a first example we can plug $C^\bullet(\fn)$ itself into the derived pure spinor functor $A^\bullet$ and study the associated multiplet.

\begin{lem}
	There is a natural equivalence
	\begin{equation*}
	A^\bullet\big(C^\bullet(\fn)\big) \simeq \Omega^\bullet(N).
	\end{equation*}
	
\end{lem}

\begin{proof}
	First, we can describe
	\begin{equation*}
	A^\bullet\big(C^\bullet(\fn)\big) = \left( C^\infty(N) \otimes C^\bullet(\fn) , \cD = \lambda^\alpha \sR(d_\alpha) + v^\mu \sR(e_\mu) + \d_{\mathrm{CE}} \right) .
	\end{equation*}
	Let us write $V$ for the even part $N_2$ of $N$, viewed as an affine space. We can identify $C^\infty(N) = C^\infty(V) \otimes \CC[\theta^\alpha]$ and $C^\bullet(\fn) = \CC[\lambda^\alpha , v^\mu]$. The differential takes the explicit form
	\begin{equation*}
	\cD = \lambda^\alpha \frac{\partial}{\partial \theta} - \lambda^\alpha f^\mu_{\alpha \beta} \theta^\beta \frac{\partial}{\partial x^\mu} + v^\mu \frac{\partial}{\partial x^\mu} + \big(\lambda^\alpha f^\mu_{\alpha \beta} \lambda^\beta\big) \frac{\partial}{\partial v^\mu} .
	\end{equation*}
	On the other hand, $N$ is parallelizable, therefore its de Rham complex takes the form
	\begin{equation*}
	\Omega^\bullet(N) \cong C^\infty(N) \otimes \Sym^\bullet(\fn_1^*) \otimes \wedge^\bullet\fn_2^* .
	\end{equation*}
	Identifying $\d\theta^\alpha = \lambda^\alpha$ and $\d x^\mu = v^\mu$, the de Rham differential takes the form
	\begin{equation*}
	\d_{\mathrm{dR}} = \lambda^\alpha \frac{\partial}{\partial \theta^\alpha} + v^\mu \frac{\partial}{\partial x^\mu} .
	\end{equation*}
	Further, $\fg$ acts on the de Rham complex via left translation making it a strict multiplet. The map defined in coordinates via
	\begin{equation*}
	\big( A^\bullet\big(C^\bullet(\fn)\big) , \cD , \sL \big) \longrightarrow \left( \Omega^\bullet(N) , \d_{\mathrm{dR}} , \sL \right), \qquad \big(x^\mu, \lambda^\alpha, \theta^\alpha , v^\mu\big) \mapsto \big(x^\mu, \lambda^\alpha, \theta^\alpha, v^\mu + \lambda f^\mu \theta\big)
	\end{equation*}
	is a quasi-isomorphism of $\fg$-multiplets. Therefore, we can identify the multiplet associated to $C^\bullet(\fn)$ itself as the differential forms on the super Lie group $N$.
\end{proof}

\begin{rmk}
	We can further compute cohomology with respect to $\cD_0 = \lambda^\alpha \frac{\partial}{\partial \theta^\alpha}$ and obtain a~retraction\footnote{Recall that a retraction between two cochain complexes is a diagram as in~\eqref{eq: dR transfer}, such that $p \circ i = \id_{\Omega^\bullet(V)}$ and $i \circ p - \id_{\Omega^\bullet(N)} = \cD_0 \circ h + h \circ \cD_0$.} to the de Rham complex on the spacetime manifold $V$:
	\begin{equation} \label{eq: dR transfer}
	\begin{tikzcd}
	\arrow[loop left]{l}{h}\big(\Omega^\bullet(N) , \cD_0\big)\arrow[r, shift left, "p"] & \big(\Omega^\bullet(V) , 0\big).\arrow[l, shift left, "i"]
	\end{tikzcd}
	\end{equation}
	Here, $i$ is the embedding of the factor of polynomial degree 0 in the $\lambda$ and $\theta$ variables, and $p$ is the obvious projection. The homotopy $h$ is given by $h=\theta^\alpha \frac{\partial}{\partial \lambda^\alpha}$.
	The induced differential via homotopy transfer is the de Rham differential on $V$. The module structure, however, is no longer strict. In fact the strict part now vanishes, but there are now quadratic pieces appearing. These take the form
	\begin{equation*}
	\rho^{(2)} (Q , Q) = \iota_{[Q, Q ]} \colon \ \Omega^k(V) \longrightarrow \Omega^{k-1}(X) .
	\end{equation*}
	Here we view $[Q, Q]$ as a constant vector field on $V$ and $\iota$ denotes the contraction of a differential form with a vector field.
	
	
	There is of course a further quasi-isomorphism of $\fg$-modules to the trivial $\fg$-module $\CC$. In this spirit, each of the multiplets $A^\bullet\big(C^\bullet(\fn)\big)$, $\Omega^\bullet(N)$, and $\Omega^\bullet(V)$ can be viewed as free resolutions of the trivial module---either free over spacetime (i.e., as $C^\infty(V)$-modules), or even free over superspace (i.e., as $C^\infty(N)$-modules). Note, however, that the trivial module does not form a $\fg$-multiplet since the translations do not act geometrically. Therefore this last quasi-isomorphism is just a quasi-isomorphism of $\fg$-modules.
\end{rmk}


\subsection{Component fields and homotopy transfer} \label{sec: components}
The multiplet $A^\bullet(\Gamma)$ given by applying the pure spinor functor does not at first glance resemble the more standard component field formulations known from physics. These component field multiplets are distinguished by the fact that they are given by dgs vector bundles whose total rank (as vector bundles over spacetime) is finite. In other words, component field multiplets are usually defined with the assumption that they only contain a finite number of component fields. With some care, it is always possible to choose a homotopy representative for $A^\bullet(\Gamma)$ that satisfies this condition, as long as $\Gamma$ satisfies a finiteness condition. In other words, we will establish the following lemma.

\begin{lem} \label{fd_model_lemma}
	Let $\Gamma$ be a $C^\bullet(\fn)$-module such that only finitely many cohomology groups $\mathrm{H}^\bullet(\Gamma)$ are non-vanishing and each of these is finitely generated as an $R$-module. Then $A^\bullet(\Gamma)$ is quasi-isomorphic to a multiplet of finite rank.
\end{lem}

In the rest of this section we will explain an algorithm that provides explicit finitely generated component field representatives for $A^\bullet(\Gamma)$, thus establishing the lemma. In~\cite{perspectives}, a general procedure to extract such a component field multiplet out of the underived model $A^\bullet_{R/I}(\Gamma)$ was explained. In essence, this is done by homotopy transfer: one identifies a retraction to another quasi-isomorphic complex and then applies homotopy transfer to the structures present for the multiplet. By construction, this yields a quasi-isomorphic and thus physically equivalent multiplet. Crucially, the component field multiplets obtained in that way are not necessarily strict anymore. Higher-order terms in the module structure can arise during the transfer.

Let us begin by summarizing the procedure for the underived pure spinor formalism $A^\bullet_{R/I}$ and then describe the generalization to $A^\bullet$.

\subsubsection[Minimal models for A\^{}bullet\_\{R/I\}]{Minimal models for $\boldsymbol{A^\bullet_{R/I}}$}

Recall that the differential on $\big(A^\bullet_{R/I} (\Gamma),\cD\big)$ admits an obvious splitting
\begin{equation*}
\cD = \cD_0 + \cD_1 = \lambda^\alpha \frac{\partial}{\partial \theta^\alpha} - \lambda^\alpha f^\mu_{\alpha \beta} \theta^\beta \frac{\partial}{\partial x^\mu} ,
\end{equation*}
which can be viewed by equipping $A^\bullet_{R/I}(\Gamma)$ with a filtration by polynomial degree on $V$ and splitting the differential into the terms that preserve and lower filtered degree.
(To filter by polynomial degree on~$V$ is imprecise, since that filtration is not complete. A complete filtration that serves the same purpose can be constructed as explained in~\cite{perspectives}, roughly speaking by assigning both $\lambda$ and $\theta$ filtered degree one. Alternatively, under the assumption that the generator of scalings is included in~$\fg_0$, one can use the weight grading as discussed in the next subsection.)
We can take cohomology with respect to $\cD_0$ and then perform homotopy transfer along the diagram
\begin{equation*}
\begin{tikzcd}
\arrow[loop left]{l}{h}(A^\bullet (\Gamma) , \cD_0)\arrow[r, shift left, "p"] &(\mathrm H^\bullet(A^\bullet(\Gamma), \cD_0) , \, 0)\arrow[l, shift left, "i"] .
\end{tikzcd}
\end{equation*}
We observe, by \cite[Theorem 10.3.15]{LodayVallette}, that although the transfer data depends on a choice of section for the projection onto the cohomology, the multiplet obtained by homotopy transfer is independent of this choice up to isomorphism. In this example, it is always ``minimal'' in the sense that its differential does not contain terms of order zero in differential operators.
Intuitively, this means that we cannot take any further cohomology without leaving the category of multiplets. (In more general examples coming from $C^\bu(\fn)$-modules, ``minimal'' is not related to having no terms of order zero in the differential; the massive Klein--Gordon field is a simple example.)

Starting from a multiplet of the form $A^\bullet_{R/I}(\Gamma)$, we will refer to the corresponding minimal multiplet as $\mu A^\bullet_{R/I}(\Gamma)$.
One can identify the $\cD_0$-cohomology with the Koszul cohomology of $\Gamma$, tensored with functions on spacetime,
\begin{equation*}
\mathrm H^\bullet(A_{R/I}^\bullet(\Gamma)) = C^\infty(V) \otimes \mathrm H^\bullet(K^\bullet(\Gamma)) .
\end{equation*}
The Koszul cohomology is conveniently computed by a minimal free resolution $L$ of $\Gamma$ in $R$-modules. The minimal multiplet takes the form
\begin{equation*}
\mu A^\bullet_{R/I}(\Gamma) = \big( C^\infty(V) \otimes (L \otimes_R \CC) , \cD' ,\rho' \big) ,
\end{equation*}
where $\cD'$ is the differential induced from $\cD_1$ and $\rho'$ the module structure induced from $\sL$ via homotopy transfer. We refer to~\cite{perspectives} for details. For $\Gamma$ a finitely generated module, Hilbert's syzygy theorem states that the minimal free resolution exists, consists of finitely generated modules and its length is less or equal than $\dim(\fn_1)$ (see for example~\cite[Theorem 1.13]{MR1322960}; for a~discussion in the equivariant case we refer to~\cite[Proposition 2.4.9, Remark 2.4.10]{Galetto}). Therefore,~$\mu A_{R/I}^\bullet(\Gamma)$ is indeed of finite rank over spacetime, i.e., it provides a reasonable component field multiplet.

\subsubsection[Minimal models for A\^{}bullet]{Minimal models for $\boldsymbol{A^\bullet}$}

Let us now discuss a generalization of this procedure to the functor $A^\bullet$. We will assume in this section that $\Gamma$ carries an action of the abelian Lie algebra $\RR$ compatible with the action of $\RR$ on $\fn$ where $\fn_i$ has $\RR$-weight $i$. This will be used to construct a filtration at the very end of this section (however, this assumption is not required for Lemma~\ref{fd_model_lemma}).

Recall that the differential on $A^\bullet(\Gamma)$ takes the form
\begin{equation*}
\cD = \lambda^\alpha \frac{\partial}{\partial \theta^\alpha} - \lambda^\alpha f^\mu_{\alpha \beta} \theta^\beta \frac{\partial}{\partial x^\mu} + \d_\Gamma + v^\mu
\frac{\partial}{\partial x^\mu} .
\end{equation*}
One can construct a finite-dimensional component field model in two steps. First one takes cohomology with respect to the internal differential of the module $\d_\Gamma$ and performs homotopy transfer along the diagram
\begin{equation*}
\begin{tikzcd}
\arrow[loop left]{l}{h}(A^\bullet (\Gamma) , \d_\Gamma)\arrow[r, shift left, "p"] &(C^\infty(N) \otimes \mathrm H^\bullet(\Gamma) ) , \, 0)\arrow[l, shift left, "i"] .
\end{tikzcd}
\end{equation*}
The differential and the $C^\bullet(\fn)$-module structure of the resulting multiplet may contain additional pieces induced from homotopy transfer. Since $\mathrm H^\bullet(\Gamma)$ is a $C^\bullet(\fn)$-module with vanishing differential, each homogeneous summand $\mathrm H^k(\Gamma)$ carries the structure of an $R/I$-module. Thus, we can proceed by taking cohomology with respect to the Koszul differential $\cD_0 = \lambda^\alpha \frac{\partial}{\partial \theta^\alpha}$, and applying the homotopy transfer along a retraction on to this cohomology. As before, the $\cD_0$-cohomology is computed by minimal free resolutions of the individual $\mathrm H^k(\Gamma)$ in $R$-modules, and as before, it is independent of the choice of transfer datum. The resulting multiplet is thus of the form
\begin{equation*}
C^\infty(V) \otimes \bigg(\bigoplus_k L^\bullet_k \otimes_R \CC\bigg) ,
\end{equation*}
where $L^\bullet_k$ is the minimal free resolution of $\mathrm H^k(\Gamma)$. As long as the cohomology $\mathrm H^\bullet(\Gamma)$ is bounded, this is already a finite rank vector bundle over spacetime. The field content of this multiplet is just the field content of the direct sum of all the minimal multiplets associated to the cohomology groups, i.e.,
\begin{equation*}
\bigoplus_k \mu A^\bu\big(\mathrm H^k(\Gamma)\big) .
\end{equation*}
At this stage we have already obtained a finite-dimensional model, as required by Lemma \ref{fd_model_lemma}. However, additional acyclic differentials induced by homotopy transfer can still be present. At this final stage, we will use the filtration on each multiplet $\mu A^\bu\big(\mathrm H^k(\Gamma)\big)$ by weight with respect to the $\RR$-action on $\Gamma$. We define $\mu A^\bullet(\Gamma)$ by taking the cohomology with respect to the summand of the total differential of filtered degree zero (in other words, we pass to the $E_1$ page of the associated spectral sequence), and apply homotopy transfer. We will illustrate this procedures using examples in~Section~\ref{sec: examples}.


\section{An equivalence of categories} \label{sec: equiv}
We now show that the derived pure spinor superfield formalism provides an equivalence of categories between the dg categories of $\fg_0$-equivariant $C^\bullet(\fn)$-modules and $\fg$-multiplets. This implies in particular that, up to quasi-isomorphism, every $\fg$-multiplet can be constructed using the derived pure spinor superfield formalism.

\begin{rmk}
	Recall that for a general Lie algebra there is a Koszul duality equivalence between the dg-categories of $C_\bullet(\fg)$-comodules and $U(\fg)$-modules. If $\fg$ is, for instance, finite dimensional, we can instead consider $C^\bullet(\fg)$-modules. For a relevant discussion in a similar context, see~\mbox{\cite[Sections~7 and~8]{CostelloYangian}}. Relatedly, Kapranov~\cite{Kapranov} gives a formulation of Koszul duality that establishes a Quillen equivalence between the categories of $D$-modules and $\Omega^\bu$-modules on the same space. Our results show that the pure spinor formalism admits a natural derived generalization that is closely related to Kapranov's construction; however, in our setting, one is working on a supermanifold, and asks for appropriate equivariance conditions. Alternatively, our procedure can be viewed in two steps, as an explicit form of commutative/Lie Koszul duality tailored to the examples in question, combined with an associated bundle construction that realizes a~$U(\fn)$-module as an $\fn$-equivariant dg vector bundle over $V = N_2$.
\end{rmk}

\subsection[The inverse functor: derived n-invariants]{The inverse functor: derived $\boldsymbol{\fn}$-invariants}

Any $\fg$-module is in particular a $\fn$-module, and we can thus take derived invariants with respect to $\fn$. For multiplets, this defines a functor in the opposite direction to the functor $A^\bullet$, assigning a strict $C^\bullet(\fn)$-module to a multiplet. The resulting $C^\bullet(\fn)$-module is also $\fg_0$-equivariant.
The functor takes the form
\begin{equation*}
C^\bullet(\fn , -) \colon \ \Mult_{\fg}^{\mathrm{strict}} \longrightarrow \Mod^{\fg_0}_{C^\bullet(\fn)},
\end{equation*}
It maps the multiplet	$(E,D,\rho)$ to $C^\bullet(\fn , \cE)$; this Chevalley--Eilenberg complex can be written more explicitly as
\begin{equation*}
C^\bullet(\fn , \cE) = \big( C^\bullet(\fn) \otimes \cE , \d_{\mathrm{CE}} + D + \lambda^\alpha \rho(d_\alpha) + v^\mu \rho(e_\mu) \big) .
\end{equation*}
It is a strict $C^\bullet(\fn)$-module; the action is on $C^\bullet(\fn)$ via multiplication and on $\cE$ via the identity. Morphisms are mapped according to the rule
\begin{equation*}
\psi \mapsto
\id_{C^\bullet(\fn)} \otimes \psi .
\end{equation*}

\begin{rmk}
	The assignment on the level of objects is also well defined for not necessarily strict modules. Then the higher-order terms of the $\fg$-module structure enter the differential such that the term $\lambda^\alpha \rho(d_\alpha)$ is replaced by
	\begin{equation*}
	\lambda \cdot \rho := \sum_{k=1}^{\infty} \lambda^{\alpha_1} \cdots \lambda^{\alpha_k} \rho^{(k)}(d_{\alpha_1} , \dots , d_{\alpha_k}).
	\end{equation*}
	Notice that the output is still a strict $C^\bullet(\fn)$-module.
\end{rmk}
There are several ways to intuitively understand the fact that the inverse functor is given by the derived invariants of supertranslations. One can of course appeal to the general structure of Koszul duality. On a more down-to-earth level, one can recall that the component fields of a supermultiplet in the usual pure spinor superfield formalism correspond to the generators of the minimal free resolution of that module over $R = C^\bu(\Pi \fn_1)$. The resolution differential was then proven in~\cite{perspectives} to agree with those supersymmetry transformations that are order zero in spacetime derivatives, verifying a conjecture of Berkovits. Since spacetime derivatives act by zero on translation-invariant sections of~$E$, it is clear that one can think of the minimal free resolution of the module as arising from the derived $\Pi \fn_1$-invariants of translation-invariant sections: $R$ consists of the Lie algebra cochains of~$\Pi \fn_1$, the generators of the free graded $R$-module arise from the component fields of the multiplet, and the differential encodes the action of supersymmetry on translation-invariant sections. It is clear that this story is just a two-step procedure to compute derived $\fn$-invariants by first taking $\fn_2$-invariants, and then accounting for the action of supersymmetry.

To flesh this story out, we now describe the module $C^\bullet(\fn , \cE)$ in some more detail and sketch how its cohomology can be computed. Recall that, since $\cE$ is a multiplet, the action of $\fn_2$ is just given by derivatives along the coordinate directions
\begin{equation*}
\rho(e_\mu) = \frac{\partial}{\partial x^\mu} .
\end{equation*}
Thus taking cohomology with respect to the term $v^\mu \rho(e_\mu)$ in the differential means restricting to translation invariant sections of $\cE$ and eliminating $v^\mu$. Denoting the fiber of $E$ over 0 by $E_0$ we find a quasi-isomorphism
\begin{equation} \label{eq: compute C^bu}
C^\bullet(\fn ,\cE) \simeq \big( R \otimes E_0^\bullet , \lambda \cdot \rho_{\mathrm{constants}} \big) .
\end{equation}
Note that $E_0^\bullet$ carries a $\ZZ \times \ZZ/2\ZZ$-grading since it comes from a dgs vector bundle. In addition, there is an integer grading by polynomial degree in $\lambda$ present. In the following, we will label cohomology groups by the former degree. Then each cohomology group is an $\fg_0$-equivariant~$R/I$-module graded by polynomial degree in $\lambda$. If the cohomology is concentrated in a single degree, the complex on the right hand-side can be viewed as the minimal free resolution of the $R/I$-module forming the cohomology. Note that this is indeed a minimal free resolution, since all terms in the differential are of nonzero polynomial degree in $\lambda$. It will follow from the theorem below that this $R/I$-module is precisely the algebraic input datum the multiplet can be constructed from in the pure spinor superfield formalism, i.e., by applying the functor $A^\bullet_{R/I}$.


\subsection{Main theorem and proof}
Let us now show that the functors $A^\bullet$ and $C^\bullet(\fn,-)$ induce an equivalence of dg-categories between the homotopy categories of multiplets and equivariant $C^\bullet({\fn})$-modules.

\begin{thm} \label{thm: equiv}
	$A^\bullet$ and $C^\bullet(\fn , - )$ provide an equivalence of dg-categories between $\Mult_\fg^{\text{\rm strict-ob}}$ and $\Mod^{\fg_0}_{C^\bullet({\fn})}$.
\end{thm}
\begin{proof}
	We first show that there is an equivalence of dg-functors
	\begin{equation*}
	\mathrm{id}_{\Mod^{\fg_0}_{C^\bullet({\fn})}} \to C^\bullet \circ A^\bullet .
	\end{equation*}
	We will naturally construct quasi-isomorphisms
	\begin{equation*}
	\Gamma \simeq C^\bu(\fn,A^\bu(\Gamma))
	\end{equation*}
	for each equivariant $C^\bu(\fn)$-module $(\Gamma,\d_\Gamma)$.
	
	We can explicitly describe $C^\bu(\fn,A^\bu(\Gamma))$ by
	\begin{gather*}
	\bigg( \CC\big[\lambda', v'\big] \otimes C^\infty(N) \otimes \Gamma ,\\
\qquad{}  \d_\Gamma + \big(\lambda' f^\mu \lambda'\big) \frac{\partial}{\partial v'^\mu} + v'^\mu \frac{\partial}{\partial x^\mu} + \lambda'^\alpha \sL(d_\alpha) + \lambda^\alpha \sR(d_\alpha) + v^\mu \sR(e_\mu) \bigg) .
	\end{gather*}
	Here we made a notational distinction between the generators of the $C^\bullet(\fn)$ in the construction (denoted by $\lambda'$ and $v'$) and the action of $C^\bullet(\fn)$ on $\Gamma$ (denoted by $\lambda$ and $v$).
	
	The differential contains a piece of the form
	\begin{equation*}
	\d_{\mathrm{dR}} = v'^\mu \frac{\partial}{\partial x^\mu} + \lambda'^\alpha \frac{\partial}{\partial \theta^\alpha} .
	\end{equation*}
	Thus, we can identify
	\begin{equation*}
	\left( C^\bu(\fn,A^\bu(\Gamma)) , v'^\mu \frac{\partial}{\partial x^\mu} + \lambda'^\alpha \frac{\partial}{\partial \theta^\alpha} \right) = \left( \Omega^\bu(N) , \d_{\mathrm{dR}} \right) \otimes \Gamma .
	\end{equation*}
	Since the de Rham complex on $N$ is acyclic, this complex is quasi-isomorphic to $\Gamma$. We can fix homotopy data
	\begin{equation*}
	\begin{tikzcd}
	\arrow[loop left]{l}{h} \left( \Omega^\bu(N) \otimes \Gamma , \d_{\mathrm{dR}} \right) \arrow[r, shift left, "p"] & (\Gamma , \, 0)\arrow[l, shift left, "i"] .
	\end{tikzcd}
	\end{equation*}
	It is easy to see that the only induced differential on the right hand side is $\d_\Gamma$, so that we obtain a quasi-isomorphism uniformly for all choices of $\Gamma$:
	\begin{equation*}
	C^\bu(\fn,A^\bu(\Gamma)) \simeq (\Gamma , \d_\Gamma) .
	\end{equation*}
	
	Now, for any morphism $f \colon \ \Gamma \to \Gamma'$ of $C^\bullet(\fn)$-modules, we can identify
	\begin{equation*}
	C^\bu(A^\bu(f)) = \id_{C^\bu(\fn)} \otimes \id_{C^\infty(N)} \otimes f.
	\end{equation*}
	So there is a commutative square
	\begin{equation*}
	\begin{tikzcd}
	\Gamma \arrow[d , "f"] \arrow[r] & C^\bullet(\fn , A^\bullet(\Gamma)) \arrow[d, "A^\bullet(f)"] \\
	\Gamma' \arrow[r] &C^\bullet(\fn , A^\bullet(\Gamma'))
	\end{tikzcd}
	\end{equation*}
	inducing an equivalence of hom complexes $\mathrm{Hom}(\Gamma, \Gamma') \to \mathrm{Hom}(C^\bullet(\fn , A^\bullet(\Gamma)), C^\bullet(\fn , A^\bullet(\Gamma')))$.
	
	
	Conversely, we will construct an equivalence of dg-functors
	\begin{equation*}
	A^\bullet \circ C^\bullet \to \mathrm{id}_{\Mult_\fg^{\text{\rm strict-ob}}} .
	\end{equation*}
	
	Let $(E,D,\rho)$ be a $\fg$-multiplet. We can describe $A^\bullet(C^\bu(\fn , \cE))$ explicitly by
	\begin{equation*}
	\left( \CC[\lambda, v] \otimes \cE \otimes C^\infty(N) , D + \lambda f^\mu \lambda \frac{\partial}{\partial v^\mu} + \lambda \cdot \rho + v^\mu \rho(e_\mu) + \lambda^\alpha \sR(d_\alpha) + v^\mu \sR(e_\mu) \right) .
	\end{equation*}
	We denote coordinates on the base of the vector bundle $E$ by $x^\mu$ and on $N_2$ by $y^\mu$. In this notation, we find
	\begin{equation*}
	v^\mu\left( \rho(e_\mu) + \sR(e_\mu) \right) = v^\mu \left( \frac{\partial}{\partial x^\mu} + \frac{\partial}{\partial y^\mu} \right) .
	\end{equation*}
	Taking cohomology with respect to this piece, we obtain a quasi-isomorphic complex of the form
	\begin{equation*}
	\big( \CC[\lambda, \theta] \otimes \cE , D + \lambda \cdot \rho + \lambda^\alpha \sR(d_\alpha) \big) .
	\end{equation*}
	The differential contains a piece corresponding to the Koszul differential
	\begin{equation*}
	\d_K = \cD_0 = \lambda^\alpha \frac{\partial}{\partial \theta^\alpha} .
	\end{equation*}
	Clearly, taking cohomology with respect to the Koszul differential we arrive at $\cE$. Concretely let us consider homotopy data
	\begin{equation*}
	\begin{tikzcd}
	\arrow[loop left]{l}{\cD_0^\dagger}(\CC[\lambda,\theta] \otimes \cE , \cD_0 )\arrow[r, shift left, "p"] &( \cE , \, 0)\arrow[l, shift left, "i"] ,
	\end{tikzcd}
	\end{equation*}
	where the homotopy is given by $\cD_0^\dagger = \theta \frac{\partial}{\partial \lambda}$, while $i$ is the obvious inclusion and $p$ evaluates at $\lambda = \theta = 0$. It is easy to see that the induced differential is just $D$. Thus, we find homotopy data
	\begin{equation} \label{eq: hdata}
	\begin{tikzcd}
	\arrow[loop left]{l}{\cD_0'^\dagger}(\CC[\lambda,\theta] \otimes \cE , d )\arrow[r, shift left, "p'"] &( \cE , \, D )\arrow[l, shift left, "i'"] ,
	\end{tikzcd}
	\end{equation}
	providing a quasi-isomorphism of cochain complexes. The induced maps are given by
	\begin{align}
	&i' = \sum_{n=0}^{\infty} \big(\cD_0^\dagger(D + \cD_1 + \lambda \cdot \rho)\big)^n \circ i,\nonumber \\
	&p' = p \circ \sum_{n=0}^{\infty} \big(\cD_0^\dagger(D + \cD_1 + \lambda \cdot \rho)\big)^n = p,\nonumber \\
	&\cD_0'^\dagger = \cD_0^\dagger \circ \sum_{n = 0}^{\infty} \big((D + \cD_1 + \lambda \cdot \rho) \cD_0^\dagger\big)^n.	\label{htpy_transfer_in_main_thm}
	\end{align}
	By construction, $p \circ \cD_0^\dagger = 0$ and thus $p' = p$. We note that these sums are all finite by degree reasons, since the left hand side in~\eqref{eq: hdata} is concentrated in finitely many degrees with respect to the grading by polynomial degree in the $\theta$-variables, and in each expression in equation~\eqref{htpy_transfer_in_main_thm}, the operator being raised to the power $n$ within the sum raises this $\theta$-degree. Therefore for $n$ sufficiently large, all terms in the sum defining our maps vanish.
	
	We have to check that this not only provides a quasi-isomorphism of cochain complexes, but of multiplets. Therefore we transfer the module structure induced by $\sL$ to the right hand side and check that it agrees with the original module structure $\rho$ of the multiplet.
	
	Explicitly, the transferred module structure can be described by
	\begin{align*}
	\rho_\sL^{(k)} (Q,\dots,Q)={}& p \sL(Q) \big(\cD_0'^\dagger \sL(Q)\big)^{k-1} i' \\
	={}&p \sL(Q) \left(\cD_0^\dagger \sum_{n = 0}^{\infty} ((D + \cD_1 + \lambda \cdot \rho) \cD_0^\dagger)^n \sL(Q)\right)^{k-1}\\
&{}\times \sum_{j=0}^{\infty} \big(\cD_0^\dagger(D + \cD_1 + \lambda \cdot \rho)\big)^j \circ i .
	\end{align*}
	Since $p$ projects onto $\lambda = \theta = 0$, only terms of order zero in $\lambda$ and $\theta$ can contribute. It is easy to see that the only such term is
	\begin{equation*}
	\rho_\sL^{(k)} (Q,\dots,Q) = p \epsilon \frac{\partial}{\partial \theta} \left(\cD_0^\dagger \epsilon \frac{\partial}{\partial \theta}\right)^{k-1} \cD_0^\dagger \lambda \cdot \rho^{(k)} i .
	\end{equation*}
	Here we expressed $Q$ in a basis, $Q = \epsilon^\alpha d_\alpha$.
	Noting that
	\begin{equation*}
	\left\{\cD_0^\dagger , \epsilon \frac{\partial}{\partial \theta}\right\} = \epsilon \frac{\partial}{\partial \lambda} ,
	\end{equation*}
	we deduce
	\begin{equation*}
	\rho_\sL^{(k)} (Q,\dots,Q) = \rho^{(k)} (Q,\dots,Q) .
	\end{equation*}
	This shows that $A^\bu(C^\bu(\fn , \cE)) \simeq (E,D,\rho)$ as multiplets.
	
	As before, for any morphism $\psi \colon \ \cE \to \cE'$ we can realize
	\begin{equation*}
	A^\bu(C^\bu(\psi)) = C^\bullet(\psi) \otimes \id_{C^\infty(N)} .
	\end{equation*}
	So there is a commutative diagram
	\begin{equation*}
	\begin{tikzcd}
	A^\bullet(C^\bullet(\fn,\cE)) \arrow[r] \arrow[d, "C^\bullet(\psi)"] & \cE \arrow[d , "\psi"] \\
	A^\bullet(C^\bullet(\fn,\cE')) \arrow[r] & \cE'
	\end{tikzcd}
	\end{equation*}
	inducing an equivalence of hom spaces, and hence we have an equivalence of dg-categories.
\end{proof}

\begin{rmk}
	The equivalence of dg-categories automatically induces an equivalence of the underlying homotopy categories. Hence, each multiplet is perturbatively equivalent to a multiplet constructed via the derived pure spinor superfield formalism.
\end{rmk}
\begin{thebibliography}{99}
\bibitem[1]{AKLL}
Alexandrov V., Krotov D., Losev A., Lysov V., On pure spinor superfield
 formalism, \href{https://doi.org/10.1088/1126-6708/2007/10/074}{\textit{J.~High Energy Phys.}} \textbf{2007} (2007), no.~10, 074,
 42~pages, \href{https://arxiv.org/abs/0705.2191}{arXiv:0705.2191}.

\bibitem[2]{Allocca}
Allocca M.P., $L$-infinity algebra representation theory, Ph.D.~Thesis,
 {N}orth Carolina State University, 2010, available at
 \url{https://www.proquest.com/docview/897930696}.

\bibitem[3]{BaulieuSusy}
Baulieu L., Bellon M., Ouvry S., Wallet J.C., Batalin--{V}ilkovisky analysis of
 supersymmetric systems, \href{https://doi.org/10.1016/0370-2693(90)90557-M}{\textit{Phys. Lett.~B}} \textbf{252} (1990), 387--394.

\bibitem[16]{CostelloYangian}
Costello K., Supersymmetric gauge theory and the {Y}angian, \href{https://arxiv.org/abs/1303.2632}{arXiv:1303.2632}.

\bibitem[19]{perspectives}
Eager R., Hahner F., Saberi I., Williams B.R., Perspectives on the pure spinor
 superfield formalism, \href{https://doi.org/10.1016/j.geomphys.2022.104626}{\textit{J.~Geom. Phys.}} \textbf{180} (2022), 104626,
 47~pages, \href{https://arxiv.org/abs/2111.01162}{arXiv:2111.01162}.

\bibitem[21]{MR1322960}
Eisenbud D., Commutative algebra: with a view toward algebraic geometry,
 \textit{Grad. Texts in Math.}, Vol.~150, \href{https://doi.org/10.1007/978-1-4612-5350-1}{Springer}, New York, 1995.

\bibitem[23]{ESW}
Elliott C., Safronov P., Williams B.R., A taxonomy of twists of supersymmetric
 {Y}ang--{M}ills theory, \href{https://doi.org/10.1007/s00029-022-00786-y}{\textit{Selecta Math. (N.S.)}} \textbf{28} (2022), 73,
 124~pages, \href{https://arxiv.org/abs/2002.10517}{arXiv:2002.10517}.

\bibitem[28]{Galetto}
Galetto F., Propagating weights of tori along free resolutions,
 \href{https://doi.org/10.1016/j.jsc.2015.05.004}{\textit{J.~Symbolic Comput.}} \textbf{74} (2016), 1--45, \href{https://arxiv.org/abs/1406.1900}{arXiv:1406.1900}.

\bibitem[38]{Kapranov}
Kapranov M., On {DG}-modules over the de~{R}ham complex and the vanishing
 cycles functor, in Algebraic {G}eometry ({C}hicago, {IL}, 1989),
 \textit{Lecture Notes in Math.}, Vol.~1479, \href{https://doi.org/10.1007/BFb0086264}{Springer}, Berlin, 1991, 57--86.

\bibitem[41]{KrizMay}
K\v{r}\'{\i}\v{z} I., May J.P., Operads, algebras, modules and motives,
 \textit{Ast\'erisque} \textbf{233} (1995), iv+145~pages.

\bibitem[42]{Lada}
Lada T., {$L_\infty$} algebra representations, \href{https://doi.org/10.1023/B:APCS.0000013809.71153.30}{\textit{Appl. Categ. Structures}}
 \textbf{12} (2004), 29--34.

\bibitem[43]{LodayVallette}
Loday J.-L., Vallette B., Algebraic operads, \textit{Grundlehren Math. Wiss.},
 Vol.~346, \href{https://doi.org/10.1007/978-3-642-30362-3}{Springer}, Heidelberg, 2012.

\end{thebibliography}
\end{document}
